%-----------------------------------
\chapter{Questionário e TCLE (exemplo)}
\label{apC_questionario}

Este apêndice exemplifica os instrumentos usados em pesquisas com participantes, como \textit{surveys} (Seção~\ref{sec_tipo_survey}) e avaliações de sistemas com usuários. Adapte o texto às exigências do CEP, quando aplicável.

\section{Termo de Consentimento Livre e Esclarecido}

Você está sendo convidado(a) a participar da pesquisa \enquote{\imprimirtitulo}, desenvolvida por \imprimirautor, estudante do curso de Bacharelado em Sistemas de Informação do IFNMG -- Campus Salinas, sob orientação de \imprimirorientador.

\textbf{Objetivo}: descreva em linguagem simples o objetivo da pesquisa.

\textbf{Procedimento}: sua participação consiste em responder a um questionário com duração aproximada de 10 minutos.

\textbf{Riscos e benefícios}: descreva os riscos (por exemplo, desconforto ao responder alguma pergunta) e os benefícios esperados.

\textbf{Sigilo}: as respostas são anônimas e serão usadas apenas para fins acadêmicos, conforme a Lei Geral de Proteção de Dados Pessoais (LGPD).

\textbf{Participação voluntária}: você pode desistir a qualquer momento, sem qualquer prejuízo.

\textbf{Contato}: informe o e-mail do pesquisador responsável.

\vspace{0.5cm}
\noindent ( \quad ) Declaro que li e concordo em participar da pesquisa.

\section{Questionário}

\textbf{Parte 1 -- Perfil do respondente}

\begin{enumerate}
	\item Idade: ( \quad ) até 20 anos \quad ( \quad ) 21 a 30 anos \quad ( \quad ) 31 a 40 anos \quad ( \quad ) mais de 40 anos
	\item Com que frequência você usa o sistema? ( \quad ) diariamente \quad ( \quad ) semanalmente \quad ( \quad ) raramente
\end{enumerate}

\textbf{Parte 2 -- Avaliação} (1 = discordo totalmente; 5 = concordo totalmente)

\begin{table}[htb]
	\centering
	\caption{Afirmações do questionário (exemplo).}
	\label{tab_questionario}
	\begin{tabularx}{\textwidth}{Xccccc}
		\toprule
		Afirmação & 1 & 2 & 3 & 4 & 5 \\
		\midrule
		O sistema é fácil de usar. & ( ) & ( ) & ( ) & ( ) & ( ) \\
		O sistema é rápido. & ( ) & ( ) & ( ) & ( ) & ( ) \\
		O sistema atende às minhas necessidades. & ( ) & ( ) & ( ) & ( ) & ( ) \\
		\bottomrule
	\end{tabularx}
	\fonte{Elaborado pelo autor.}
\end{table}

\textbf{Parte 3 -- Questão aberta}

\begin{enumerate}
	\setcounter{enumi}{2}
	\item Que melhorias você sugere para o sistema?
\end{enumerate}
